LLM-based agents have become increasingly capable of autonomous tool interaction through protocols like MCP~\cite{wang2026mcptox}. 
However, this capability introduces new attack surfaces: tool poisoning attacks can embed malicious instructions in tool descriptions, parameter schemas, or response payloads, causing agents to exfiltrate sensitive data or execute unauthorized commands.

Recent work has proposed various defense mechanisms, including protocol-level attestation (AttestMCP~\cite{maloyan2026breaking}), reasoning trace verification (RTV), and hybrid approaches (ReasoningGuard). 
While these defenses show promising results on standard benchmarks, their evaluation relies on \textit{static} attack scenarios that were not designed to be adversarial. 
This creates a fundamental blind spot: \textit{can these defenses be bypassed by semantically-equivalent but structurally different attacks?}

We address this gap with \textbf{SCARF}, a semantic constraint-aware red-teaming framework that:

\begin{itemize}
    \item \textbf{Automatically generates} adversarial MCP attack variants through semantic-preserving random search, maintaining functional equivalence while probing defense decision boundaries.
    \item \textbf{Provides four competitive baselines} (static, paraphrase, Best-of-N, random) for comprehensive defense evaluation.
    \item \textbf{Demonstrates} that even simple semantic perturbations can achieve 9.3\% defense evasion rate against state-of-the-art defenses, substantially higher than existing approaches.
\end{itemize}

Our key insight is that \textit{semantic preserves}---maintaining the original attack intent while varying surface-level tokens---can expose weaknesses in defense mechanisms that operate primarily on pattern matching or capability attestation.